\subsection{Digital embryos}

A \emph{digital embryo} is a computational reconstruction of a developmental
process that is close enough to the underlying biology to be wrong in
interesting ways. The promise of the field is that cheap computation can
replace some expensive experimentation: simulate the morphogen, the network,
or the tissue, and reserve the bench for questions the simulation cannot
answer. The peril is equally well known --- a simulation unconstrained by real
measurements will happily reproduce whatever its author expected. This report
takes the second half of that sentence seriously. Three classical model
systems of developmental patterning are rebuilt computationally, and every
quantitative claim is tied to real, accession-level public data: published
per-embryo morphogen measurements, a published gene-regulatory network model,
and published quantitative benchmarks. Where our computation contradicts an
initial expectation --- including our own earlier, stronger claims --- we say
so in the body of the report, not in a footnote.

The three studies correspond to the three great patterning problems of the
fly embryo, at three scales:

\begin{enumerate}
\item \textbf{Turing pattern selection (Study A).} Reaction--diffusion
systems of the type proposed by Turing \cite{turing1952} and instantiated by
Gierer and Meinhardt \cite{gierer1972} amplify a band of spatial wavelengths
out of noise. Linear theory predicts which wavenumber grows fastest; what the
\emph{saturated, nonlinear} pattern actually selects is a subtler question.
We measure the selected wavenumber across a parameter sweep and ask whether
the deviation from linear theory is itself predictable.
\item \textbf{Gene-network attractors (Study B).} The segment-polarity
network of \emph{Drosophila} is one of the best-studied Boolean models in
systems biology \cite{vondassow2000,albert2003}. We take the published
14-node Boolean model (Cell Collective id-021 \cite{helikar2012}) as ground
truth, enumerate its attractors exactly under all eight external signalling
conditions, measure the size of the wild-type basin of attraction, and ask
how much data a learning algorithm needs to recover the basin boundary.
\item \textbf{Morphogen positional information (Study C).} The Bicoid
morphogen gradient of the early fly embryo is the canonical example of
positional information \cite{wolpert1969,gregor2007}. Using the 676
per-embryo gradients and 2{,}180 immunofluorescence embryo records published
by Liu, Morrison and Gregor \cite{liu2013}, we (i) reproduce their gradient
length-scale measurements, (ii) test how the length scale depends on maternal
gene dosage, (iii) attack their threshold model of cephalic-furrow
positioning with a learned decoder, and (iv) measure how much positional
information gap-gene expression patterns carry under different dosages.
\end{enumerate}

\subsection{Why these three}

The three problems are independent in method but connected in spirit: each
asks where spatial order comes from, and each has a quantitative published
reference point that a computation can fail to meet. Study~A's reference is
the linear dispersion relation itself --- a prediction made by mathematics,
not by a measurement, which makes it the cleanest possible benchmark to break.
Study~B's reference is a published, curated, executable model: if our
attractor enumeration disagreed with it, the disagreement would be
diagnosable down to a single update rule. Study~C's references are the
hardest kind: numbers measured on real embryos, with real noise, by another
laboratory \cite{liu2013,petkova2019,gregor2007}.

\subsection{Summary of results}

Every number quoted here is computed by committed, tested code from committed
data, and re-computed when this report is built; the repository records the
exact pipeline. Section references point to full derivations and tables.

\begin{itemize}
\item \textbf{Study A (Sections \ref{sec:studyA}).} Across a 14-point sweep
of the activator decay rate, the saturated Gierer--Meinhardt pattern
systematically selects a wavenumber \emph{above} the linear fastest-growing
mode, by $14.7\%$ to $58.9\%$ (median $31\%$), with the largest relative
deviations at intermediate distance from onset --- not at threshold, where
amplitude-equation reasoning would have placed them. The deviation is real
(measured at steady state, verified against two lattice sizes) but it is
\emph{not} predictable from the linear dispersion curve by a convolutional
surrogate: a promising $n = 58$ pilot result ($+17.7\%$ held-out improvement)
\emph{failed to replicate} at $n = 171$ (CNN relative RMSE $0.292$--$0.295$
vs.\ $0.272$ for the linear formula alone). We retract the small-sample claim
and keep the phenomenon.
\item \textbf{Study B (Section \ref{sec:studyB}).} Exact attractor
enumeration of the published segment-polarity network under all eight
combinations of the external signals \textsc{hh}, \textsc{WG}, \textsc{SLP}
reproduces the published condition-dependent bistability exactly
\cite{albert2003}. The wild-type (\textsc{wg}-ON) basin occupies $0.77\%$ of
the synchronous state space under the wild-type signalling condition. The
basin boundary is learnable to $> 98\%$ test accuracy from only $\sim 60$
labelled states (of $2^{14}$), and a graph neural network matches or beats a
multilayer perceptron at every sample size; under naive uniform sampling the
classification task is degenerate (majority class accuracy $99.3\%$), which
we report as a cautionary result.
\item \textbf{Study C (Sections \ref{sec:studyC}--\ref{sec:gapdecode}).}
From 676 real Bcd-GFP gradients: the population length scale reproduces the
published value under the authors' fitting window ($\lambda_{\mathrm{median}}
= 0.189$ vs.\ published $0.165 \pm 0.007$ egg lengths \cite{liu2013}, with
our value on their exact window; the residual difference is discussed
honestly). The length scale \emph{increases with maternal bcd dosage}
(dosage-coupled length scale, DCLS): Spearman $\rho = 0.356$
($p = 7.9 \times 10^{-16}$) on anterior-weighted fitting windows, robust to
parametric refitting (OLS slope $+0.0117$ per dose, $95\%$ CI
$[0.0060, 0.0174]$, $p = 5.9 \times 10^{-5}$), and \emph{absent} on the
posterior-heavy window used in the original publication ($\rho = 0.071$,
$p = 0.12$) --- a window-sensitivity finding, not a contradiction. A
one-dimensional convolutional decoder predicts cephalic-furrow position from
the raw Bcd profile at $2.38\%$ egg length RMSE (held-out fly lines,
$R^2 = 0.854$), beating the published threshold model at $4.94\%$ EL by a
paired-bootstrap margin of $2.57\%$ EL, CI$_{95}$ $[1.99, 3.11]$
($p < 10^{-4}$). Finally, positional decodability from joint Hb+Eve
expression improves with dosage: median error $6.21\%$ EL at $3\times$ bcd
vs.\ $9.91\%$ EL at $1\times$, bootstrap CI $[3.00, 4.40]\%$ EL for the
difference.
\end{itemize}

\subsection{What counts as evidence here}

This project operates under a simple reporting rule: a number is quoted only
if it exists in a committed results file produced by committed, tested code;
results that fail to replicate are retracted in the same document that
reports them; and the negative results (Section~\ref{sec:negatives}) are
given the same space and the same figures as the positive ones. We regard the
retracted CNN-surrogate claim of Study~A as the single most valuable output
of the project, because it is the kind of result that a less careful pipeline
would have published.

\subsection{Background}

\textbf{Turing patterns.} Turing's 1952 observation \cite{turing1952} ---
that diffusion, usually a homogenising force, can destabilise a homogeneous
state when one morphogen activates and another, faster-diffusing one
inhibits --- is the founding result of mathematical developmental biology.
The linear theory is textbook material; the nonlinear selection problem it
leaves open (which mode, of the whole unstable band, does the saturated
pattern actually take?) has a large literature for specific kinetics and
geometries but no general closed answer. Study~A measures the gap between
linear prediction and nonlinear outcome on the canonical activator--inhibitor
kinetics of Gierer and Meinhardt \cite{gierer1972}.

\textbf{The segment-polarity network.} The repeated unit of the embryonic
parasegment is maintained by a small network of cross-cell signalling
(wingless, hedgehog) and intracellular regulation (engrailed, cubitus
interruptus, patched). von Dassow et al.\ showed by exhaustive ODE sampling
that the module is extraordinarily robust to kinetic parameters
\cite{vondassow2000}; Albert and Othmer then showed that a Boolean skeleton
of the same network already contains the essential logic \cite{albert2003}.
The curated Boolean version we use \cite{helikar2012} is therefore a model
with an unusually strong claim to correctness --- which is exactly why exact
numbers about it (basin sizes, learnability) are worth having.

\textbf{Positional information in the fly embryo.} Wolpert's positional
information framework \cite{wolpert1969} proposes that cells read their
position from morphogen concentration. The Bicoid gradient is its best
quantitative instance: an approximately exponential anterior morphogen
gradient \cite{gregor2007} that is read out by gap genes with a precision
approaching the physical limit set by expression noise
\cite{dubuis2013,petkova2019}. Liu, Morrison and Gregor added the crucial
dosage series --- embryos with 1$\times$ to 4$\times$ maternal \emph{bcd} ---
and showed the readout machinery tracks the input dynamically
\cite{liu2013}. Study~C is built entirely on their public data.

\subsection{Roadmap}

Section~\ref{sec:math-turing}ff.\ derives every formula used later, with
symbolic cross-checks. Sections \ref{sec:studyA}, \ref{sec:studyB} and
\ref{sec:studyC} present the three studies, each self-contained in method
and result. Section~\ref{sec:repro} documents the software, including the
\texttt{embryosim} command-line tool. Section~\ref{sec:discussion} discusses
scope and open work; Section~\ref{sec:negatives} collects the negative
results. Appendices carry the full tool and dataset manifests
(Appendix~\ref{app:manifest}), the complete attractor census
(Appendix~\ref{app:attractors}), reproduction commands
(Appendix~\ref{app:repro}), the results-file inventory
(Appendix~\ref{app:results}) and the test-suite inventory
(Appendix~\ref{app:tests}).
